\section{Two supplementary scale intervals}\label{sec:small-extensions}

We prove the two normalized source estimates used at the joins in
Section~\ref{sec:small}. Throughout this section we use the strict-hole
configuration of Lemma~\ref{pr:edge-reduction}, with an actual simple
rightmost zero $b$. Assume $N=n-1\ge N_0=6250$, put $n_0=6251$,
and write
\[
 F(u)=p(i+\rho u),\qquad \beta=X-i,\qquad
 0<\alpha=a/\rho\le X\le1.
\]
Thus $\rho=1-\Im b$ and $X=\Re b/\rho$. Use
$S=F''(\beta)/F'(\beta)$, $h=|S|$, $Q=\sum_v|\beta-v|^{-2}$,
$d=\min_v|\beta-v|$, $D=dh$, and $y_0=Q/h$, where $v$ runs over
all zeros of $F'$, counted with multiplicity.
The physical quantities $H_{\rm phys}=|p''(b)/p'(b)|$ and
$d_{\rm phys}=\min_{p'(\xi)=0}|b-\xi|$ satisfy
\[
 h=\rho H_{\rm phys},\qquad d_{\rm phys}=\rho d.
\]
Every critical label $v=r-it$ satisfies
\begin{equation}\label{sm:extension-domain}
 r<\alpha/2,\qquad \rho|v|^2/2\le t<\rho|v-\alpha|^2/2.
\end{equation}
The inset is $I(u)=-\Im u-\rho|u|^2/2$. In particular,
$I(\beta)=1-\rho(1+X^2)/2\ge1-\rho$.
The finite polar quantities $W,C,E,\theta$ are those in
\eqref{or:source-definitions}; their proof applies for every
$0<\rho<1$ used below.

\subsection{A closed-cell source criterion}

Suppose
\[
 \rho\in[p,r],\quad y_0\in[l,u],\quad
 X\in[X_-,X_+]\subseteq[0,1],\quad d>1,\quad
 h\ge H,\quad |\Im S|<J.
\]
Here $0<p\le r<1$, $0\le l\le u$, and $H,J>0$.
The radial identities \eqref{or:radial} and
\eqref{or:inset-critical}, with $T=1$, give
\begin{gather}
 \frac{N\rho}{2}<\rho h+J+y_0h,\label{sm:extension-radial-lower}\\
 (1-\rho/2)Q<\frac{N\rho}{2}+J.\label{sm:extension-radial-upper}
\end{gather}
For the second inequality the negative critical-inset term is retained:
\[
 E_c<\rho X\Re S-\rho X^2Q/2.
\]
Its substitution cancels the $X^2$ terms in $I(\beta)Q$.
Consequently, putting
\[
 B_r=\frac{r/2+J/N_0}{1-r/2},
\]
one may take any positive rational $U$ satisfying at least one of
\begin{equation}\label{sm:extension-U}
 U\ge u,\qquad U^2\ge B_r,\qquad Ul\ge B_r\quad(l>0)
\end{equation}
as an upper bound for $h/n$. The first alternative follows from
$h^2\le NQ$; the second uses the same inequality and
\eqref{sm:extension-radial-upper}; the third uses
$(h/n)y_0=Q/n<B_r$.

Choose $\delta>0$ with $r^2(1+\delta)^2/2<1$ and test at
$s=X+\delta$. Define
\begin{equation}\label{sm:extension-kappa}
 \kappa_\delta(X)=\frac{X^2}{2}
 -\frac{X\max(0,1-X\delta-\delta^2)^2}
           {4(1+\delta^2)(X/2+\delta)}.
\end{equation}
Assume the whole-cell inequalities
\begin{equation}\label{sm:extension-A}
 [r(1+\delta)/2]^2<\delta,\qquad
 X-\kappa_\delta(X)y_0<\overline A
\end{equation}
with $\overline A>0$. The first places the inverse Cayley test
strictly to the right of all original zeros, because
$s/(1+(\rho s/2)^2)>X$ follows from
$\delta>X(\rho s/2)^2$. The second and
Lemma~\ref{or:critical-comparison} give
$C_{\rm real}<\rho h\overline A$.
Set
\[
 f=\frac1{1-r^2(1+\delta)^2/2},\quad
 D_p=1+[r(1+\delta)/2]^2,\quad
 B=2D_p(f\overline A U+1/n_0).
\]
If $0<B<1$, the finite polar inequality yields
\begin{equation}\label{sm:extension-HB-upper}
 \frac Eh\le\frac Wh<
       \mathcal U:=\frac{r(f\overline A+1/H)}{1-B}.
\end{equation}
Both finite corrections $1/H$ and $1/n_0$ occur in this expression.

For a lower bound for $E/h$, choose a rational radius
$0<R<1$, a number $0<\nu<1$, and a positive integer $W_0$.
Put
\begin{gather}
 \lambda=\frac{uR}{1-R},\qquad
 e_R=u\left\{\sum_{k=2}^{16}\frac{R^k}{k}
                    +\frac{R^{17}}{17(1-R)}\right\},\qquad
 \widehat R=R-e_R,\label{sm:extension-inverse}\\
 B_R=\frac{u}{(1-R)^2(1-\lambda)^3},\qquad
 \epsilon=\frac{1001B_R}{1000H},\qquad
 g=\frac{1001}{1000}\nu B_R\widehat R,\qquad
 p_0=\frac1{1+\lambda},\quad a_0=\frac1{1-\lambda}.
 \notag
\end{gather}
The complete logarithmic tail appears in $e_R$.
For $W_0\in\{1,2,3,4,5,6,8,10\}$ use $e_0=1/d_{W_0}$, where
\[
\begin{array}{c|rrrrrrrr}
 W_0&1&2&3&4&5&6&8&10\\ \hline
 d_{W_0}&2&7&20&50&148&403&2900&22000 .
\end{array}
\]
In each case $\sum_{k=0}^{69}W_0^k/k!>d_{W_0}$, so
$e_0>e^{-W_0}$. Require
\begin{gather}
 0<R<1,\quad 0<\lambda<1,\quad \widehat R>0,\quad
 W_0^2<(1-\nu^2)(H\widehat R)^2,\label{sm:extension-gates-a}\\
 \epsilon(1+e_0)+ge_0<1-a_0e_0,\qquad
 \epsilon(1+e_0)+ge_0^2<p_0-e_0,\label{sm:extension-gates-b}\\
 \epsilon+ge_0<p_0,\qquad
 0<p_0<1<a_0,\qquad \nu H\widehat R>22/7.
 \label{sm:extension-gates-c}
\end{gather}
These are the entire-contour hypotheses
\eqref{or:contour-gates}--\eqref{or:odd-gate}.
The inverse-log and amplitude comparisons therefore give an
actual original-root count
$L_R>h\nu\widehat R/\pi-2$.
The phase bound confines these roots to
\begin{equation}\label{sm:extension-width}
 |u-\beta|<R,\qquad
 |\Re(u-\beta)|\le
 \ell_R:=\min(R,W_0/H+J\nu\widehat R/H+e_R).
\end{equation}
This is a bound on the entire inverse rectangle.

Use increasing radii $R_j$, round the slopes
$\nu_j(7/22)\widehat R_j$ down and the widths up to the
$10^{-12}$ mesh, and replace the widths by running maxima.
Write the resulting slopes and widths as $s_j,\ell_j$.
They are checked to satisfy
$0<s_1<\cdots<s_k$, $R_1<\cdots<R_k<1$, and
$0\le\ell_1\le\cdots\le\ell_k\le R_k$, with each
$\ell_j\le R_j$.
Rightmostness and $R_j<1$ give $|\Re u|\le1$ for the counted roots.
Their weight in $E(s)$ is at least
\begin{equation}\label{sm:extension-weight}
 \omega_j=\max\left(0,\min_{t\in\{1-R_j,\,1+R_j\}}
       \frac{t-r(1+t^2)/2}{(\delta+\ell_j)^2+t^2}\right).
\end{equation}
This controls the whole depth interval: the derivative numerator
of the displayed quotient is
$A-t^2+r(1-A)t$, $A=(\delta+\ell_j)^2>0$.
It has one positive root, which is a maximum. If an endpoint
lower bound is negative, the true inset is still nonnegative.
Round $\omega_j$ down to the same mesh and put
$w_j=\max_{m\ge j}\omega_m$.
The nested regions make these nonincreasing weights valid.

The actual source $\beta$ has weight at least
$b_0=(1-r)/(1+\delta^2)$. Formula \eqref{or:Abel-source} gives
\begin{equation}\label{sm:extension-Abel}
 \frac Eh>
 \mathcal L:=\sum_{j=1}^kw_j(s_j-s_{j-1})
                   +\frac{\min(0,b_0-3w_1)}{H},
 \qquad s_0=0.
\end{equation}
The source is removed exactly once from each cumulative count.
The finite loss three is consequently retained. If the source
compensation is positive, discarding it remains valid without
an upper bound on $h$. In particular, the first count lower
bound need not itself be positive.

\begin{lemma}\label{sm:extension-cell}
Under the preceding whole-cell hypotheses, a collection satisfying
\eqref{sm:extension-gates-a}--\eqref{sm:extension-gates-c},
the stated nesting and rounding conditions, and
$\mathcal L>\mathcal U$, excludes the entire closed cell.
\end{lemma}
\begin{proof}
Equations \eqref{sm:extension-HB-upper} and
\eqref{sm:extension-Abel} are contradictory.
\end{proof}

Equation \eqref{sm:extension-A} has a polynomial verification.
The dependence on $y_0$ is affine, so it suffices to check
$y_0=l,u$. On the active interval
$X<(1-\delta^2)/\delta$,
\begin{equation}\label{sm:extension-kappa-active}
 \kappa_\delta(X)=
 \frac{X[X-(1-\delta)^2][X+(1+\delta)^2]}
             {2(1+\delta^2)(X+2\delta)}.
\end{equation}
On the other interval it is $X^2/2$.
Clear the positive denominators in
$\overline A-X+\kappa_\delta(X)l$ and
$\overline A-X+\kappa_\delta(X)u$ separately on those two
closed pieces. The rational Bernstein rule in
Section~\ref{sec:certificates}, with exact dyadic subdivisions,
proves positivity. The two closed pieces include their common
endpoint, and both $y_0$ endpoints are checked for either sign
of $\kappa_\delta$.

\subsection{The moment interval up to \texorpdfstring{$\rho=39/500$}{rho=39/500}}

\begin{lemma}\label{sm:T-moment-extension}
There is no strict-hole configuration in the normalized
range
\[
 \frac3{50}\le\rho\le\frac{39}{500},\qquad
 0<\alpha\le X\le1,\qquad \frac54\le y_0\le\frac94
\]
satisfying
\begin{equation}\label{sm:moment-extension-input}
 D\ge139,\quad Qd/h<9/4,\quad |\Im S|<4y_0,\quad
 H_{\rm phys}>
 \frac{N-(2637/32)/d_{\rm phys}}{2(1+a)}.
\end{equation}
\end{lemma}
\begin{proof}
The source criterion requires $d>1$. The lower intersection of the two
critical-boundary circles at $r=\alpha/2$ has depth
\[
 \eta=\frac{1-\sqrt{1-\rho^2\alpha^2/4}}{\rho}
      <\frac{\rho\alpha^2}{4}.
\]
The closest-point comparison gives
$d^2>(X-\alpha/2)^2+(1-\eta)^2$.
To see why it controls the whole crescent, on the lower hole arc put
$u=r-\alpha\le-\alpha/2$. The relevant scalar product is
\[
 (X-\alpha)u+(1/\rho-1)\sqrt{\rho^{-2}-u^2}.
\]
Its derivative is positive for $u\le0$. The largest allowed real
coordinate therefore gives the nearest permitted arc point.
The upper arc and the outward radial parts are farther away.
Consequently
\[
 d^2>1+\frac{X^2}{4}-\frac{\rho\alpha^2}{2}
 \ge1+(1/4-\rho/2)X^2>1,
\]
since $1/4-(39/500)/2>0$.

Using $d>1$, $a=\rho\alpha\le\rho$, and $N\rho>2637/32$,
the last hypothesis in \eqref{sm:moment-extension-input} yields
\begin{equation}\label{sm:moment-extension-field}
 h>\frac{N\rho-2637/32}{2(1+\rho)}.
\end{equation}
The bound $D\ge139$ is an input hypothesis; the table below gives
lower bounds $h\ge H$, starting with $H=137$.

Use the five scale intervals
\[
 [.06,.064],\ [.064,.068],\ [.068,.072],\
 [.072,.075],\ [.075,.078]
\]
and the four moment intervals
\[
 Y_1=[5/4,3/2],\quad Y_2=[3/2,7/4],\quad
 Y_3=[7/4,2],\quad Y_4=[2,9/4].
\]
Every pair is a closed cell. On the $j$th moment interval
put $J=4u$ and use the following exact parameters:
\begin{center}
\begin{tabular}{c|rrrr}
\toprule
$j$&1&2&3&4\\\midrule
$\delta$&$1/2$&$11/20$&$3/5$&$13/20$\\
$10^6\overline A$&428583&358161&304796&262971\\
$J$&6&7&8&9\\
\bottomrule
\end{tabular}
\end{center}
Take $X_-=0$, $X_+=1$ in every cell. The lower fields and upper
ratios $U$ are:
\begin{center}\small
\begin{tabular}{c|r|rrrr}
\toprule
$[\rho_-,\rho_+]$&$H$&
$10^{12}U_1$&$10^{12}U_2$&$10^{12}U_3$&$10^{12}U_4$\\
\midrule
$[.060,.064]$&137&27239669422&22809917356&19645808737&17272727273\\
$[.064,.068]$&148&28952380953&24237405107&20869565218&18343685301\\
$[.068,.072]$&159&30672199171&25670816045&22098399526&19419087137\\
$[.072,.075]$&170&31966753247&26749783550&23023376624&20228571429\\
$[.075,.078]$&179&33265348596&27832119321&23951241267&21040582727\\
\bottomrule
\end{tabular}
\end{center}
For each cell $[p,r]\times[l,u]$, its $H$ is at most one of
\[
 \frac{N_0p-2637/32}{2(1+r)},\qquad
 \frac{N_0p/2-J}{r+u}.
\]
The first comparison uses \eqref{sm:moment-extension-field}
and the second uses \eqref{sm:extension-radial-lower}.
Thus $h\ge H$ in each cell.
The displayed $U$ satisfy \eqref{sm:extension-U}.
Every polynomial comparison in \eqref{sm:extension-A} is
verified by \eqref{sm:extension-kappa-active}.

The fixed contour lists have, in the order of this table,
\[
\begin{array}{rrrr}
55&49&42&37\\
56&50&44&39\\
56&50&45&40\\
58&50&45&40\\
58&51&47&42
\end{array}
\]
members, for a total of $954$.
Their exact triples $(R,\nu,W_0)$ are the rows with family
\textup{M} in
\cert{small_extensions/small_extensions_contours.csv};
the indices begin at zero, with the scale index preceding the
moment index. They are also recorded in
\cert{TMOMENT6251_PROPOSAL.json}.
The exact checker \cert{certify_Tmoment6251_complete.py}
reconstructs \eqref{sm:extension-inverse}--\eqref{sm:extension-Abel},
including every complete contour and both polynomial pieces.
Every cell satisfies
\[
 \mathcal L-\mathcal U>\frac{14}{10000}>0.
\]
Lemma~\ref{sm:extension-cell} excludes the twenty cells.

The scale and moment intervals include all their common endpoints.
For $N\ge N_0$, both lower fields increase, while $J/N$ and
$1/n$ in the upper bounds decrease. All contour, polynomial,
and source-weight bounds are fixed on their closed cells.
The contradiction therefore holds on the entire degree ray.
\end{proof}

\subsection{The high-\texorpdfstring{$X$}{X} interval up to
\texorpdfstring{$\rho=31/500$}{rho=31/500}}

\begin{lemma}\label{sm:T-highX-extension}
There is no strict-hole configuration with
\[
 \rho_L:=\frac{30501}{500000}\le\rho\le
 \rho_U:=\frac{31}{500},\qquad
 X_0:=\frac{69}{100}\le X\le1,\qquad
 0<\alpha\le X .
\]
\end{lemma}
\begin{proof}
The distance calculation preceding \eqref{or:gap-T} remains
valid on this interval. The endpoint inequality
\[
 \frac2{125}-\rho_U\left(\frac14+\frac1{15625}\right)>0
\]
gives $\eta<\alpha^2/125$ on the whole interval.
The nearest-point comparison above gives
\[
 d^2>(X-\alpha/2)^2+(1-\alpha^2/125)^2
 \ge1+\frac{117}{500}X^2=:G(X).
\]
Write $G_0=G(X_0)$.
Lemma~\ref{sm:normalized-D25} applies because
$1/100<\rho_L\le\rho\le\rho_U<9/100$; it gives $D>25$.
Proposition~\ref{mf:D25} therefore gives
\[
 y_0<3/d<3/\sqrt{G_0}<57/20.
\]
The unphased radial identity gives
\[
 h>\frac{N_0\rho_L}{2(1+57/20)}>49,\qquad
 D^2>49^2G_0>51^2.
\]
In the matrix envelope \eqref{mf:coarse-envelope}, put
$S/h=c+is$ with $c\ge0$.
If $|s|\ge7/10$, then $c<18/25$ and its right side is at most
\[
 \frac{21/5}{51^2}
 +\frac{(11/2)(18/25)}{51}
 +\frac{1-(7/10)^2}{4}-\frac{(7/10)^2}{2}<0.
\]
Thus $|\Im S|<(7/10)h$ and $\mu=1/D<1/50$.

\smallskip
\noindent\emph{A joint source and moment estimate.}
Take $\delta=3/5$ and set
\[
 b_0=\frac{1-\rho_U}{1+\delta^2},\quad
 f=[1-\rho_U^2(1+\delta)^2/2]^{-1},\quad
 D_p=1+[\rho_U(1+\delta)/2]^2.
\]
The source $\beta$ and the finite polar inequality give
\begin{equation}\label{sm:highX-K}
 \rho h A_\delta>K,\qquad
 A_\delta=X-\kappa_\delta(X)y_0,\qquad
 K=\min_{\rho\in\{\rho_L,\rho_U\}}
 \frac{n_0\rho b_0/(n_0\rho+2D_pb_0)-\rho}{f}.
\end{equation}
The expression minimized is concave in $\rho$ and increases
with $n$, as in \eqref{or:direct-K}. Exact evaluation gives
\[
 K=\frac{1793507148275409343}{2882630131640625000}
       >\frac{31}{50}.
\]
Hence $A_\delta>0$.
On the entire $X$ interval,
\begin{equation}\label{sm:highX-kappa}
 \kappa_\delta(X)>c_0X^2,\qquad c_0=12/25.
\end{equation}
The active endpoint $(1-\delta^2)/\delta=16/15$ is beyond one,
so this is a single rational-polynomial comparison.

Put $v=Xy_0$. Then $0<v<1/c_0$ and
$A_\delta<X(1-c_0v)$.
The phase bound in the radial identity gives
$N\rho/2<hR$, where
\[
 R:=\beta_0+I_0y_0,\qquad
 \beta_0=7/10+(3/10)\rho_U,\qquad
 I_0=1-\rho_L(1+X_0^2)/2.
\]
Multiplying this estimate and \eqref{sm:highX-K} for the same
configuration gives, for $q=Q/h^2=y_0/h$,
\begin{equation}\label{sm:highX-joint}
 q^2<\frac{2y_0^2A_\delta R}{NK}
 \le\frac{2v^2(1-c_0v)
            [\beta_0/X_0+I_0v/X_0^2]}{N_0K}.
\end{equation}
The second step uses the nonnegative factor $1-c_0v$ and
the decrease of $1/X,1/X^2$ on $X\ge X_0$.
The polynomial
\begin{equation}\label{sm:highX-q-polynomial}
 N_0K-800v^2(1-c_0v)
             [\beta_0/X_0+I_0v/X_0^2]
\end{equation}
is positive throughout $[0,1/c_0]$.
Clearing its positive denominator and using exact Bernstein
subdivision proves this sign, together with
\eqref{sm:highX-kappa}. Thus $q<1/20$.
Proposition~\ref{lt:q05mu02} applies with $q<1/20$ and $\mu<1/50$.
It gives two distinct original zeros, different
from $\beta$, of the form $u_\pm=\beta+w_\pm/S$, with
\begin{equation}\label{sm:highX-pair}
 |w_\pm|<17/2,\qquad
 |(w_++w_-)/2|<19/10,\qquad |\Im S|<7y_0.
\end{equation}

\smallskip
\noindent\emph{The entire range $h\ge95$.}
Suppose $h\ge95$. Then $D^2>95^2G_0>10000$.
Proposition~\ref{mf:D100}, followed by the wide row of
Proposition~\ref{lt:wideD100}, gives
\[
 Qd/h<9/4,\quad q<9/400,\quad \mu<1/100,\quad
 |\Im S|<(23/5)y_0.
\]
In particular $y_0^2G_0<81/16$ and $y_0<427/200$.
We apply Lemma~\ref{sm:extension-cell} on the ten moment
intervals in the following table, taking the entire scale
interval $[\rho_L,\rho_U]$ in each row and $X\in[X_0,1]$.
Here $J=(23/5)u$, where $u$ is the upper moment endpoint.
\begin{center}\small
\begin{tabular}{c|rcrrr}
\toprule
$[l,u]$&$H$&$\delta$&$10^6\overline A$&$10^{12}U$&contours\\
\midrule
$[0,1/4]$&607&$1/200$&1000001&179392392668&149\\
$[1/4,1/2]$&335&$1/5$&924589&129486068112&117\\
$[1/2,3/4]$&230&$3/10$&803342&65122807018&95\\
$[3/4,1]$&175&$7/20$&679683&43668386653&78\\
$[1,5/4]$&140&$9/20$&528001&32941176471&64\\
$[5/4,3/2]$&117&$1/2$&428583&26504850362&52\\
$[3/2,7/4]$&100&$11/20$&355372&22213966289&44\\
$[7/4,15/8]$&95&$3/5$&285411&19094795814&39\\
$[15/8,2]$&95&$3/5$&256512&17872445821&37\\
$[2,427/200]$&95&$3/5$&227612&16806687307&34\\
\bottomrule
\end{tabular}
\end{center}
Each $H$ is either $95$ or at most
$(N_0\rho_L/2-J)/(\rho_U+u)$. Thus $h\ge H$ follows from
the hypothesis $h\ge95$ and the radial identity.
The $U$ satisfy \eqref{sm:extension-U}.
At both moment endpoints, the entire source polynomials
are checked on $[X_0,1]$, using the exact active/inactive
split in \eqref{sm:extension-kappa-active}.

The $709$ exact triples are the family \textup{X} rows of
\cert{small_extensions/small_extensions_contours.csv},
with the cells in the printed order; they are also in
\cert{T6251_HIGHX_PROPOSAL.json}.
The checker \cert{certify_T6251_highX_complete.py}
reconstructs all the predicates in
Lemma~\ref{sm:extension-cell}. In every one of the ten cells,
\[
 \mathcal L-\mathcal U>\frac6{10000}>0.
\]
These ten cells exclude the entire range $h\ge95$.

\smallskip
\noindent\emph{The pair when $49<h<95$.}
It remains to use the pair \eqref{sm:highX-pair}.
For $\delta=3/5$,
\[
 y_0A_\delta<v-c_0v^2\le\frac1{4c_0}=\frac{25}{48}.
\]
The radial identity and $|\Im S|<7y_0$ give
$\rho(N-2h)<2y_0(h+7)$.
Since $A_\delta>0$ and $N-2h>0$,
\[
 \rho h A_\delta
 <\frac{2(25/48)95\cdot102}{6250-190}.
\]
Let $B_{95}=2D_p(95f+1)/6251<1$.
Here $A_\delta<X\le1$ by \eqref{sm:highX-kappa}, so this
is a valid bound for the complete polar denominator.
Lemma~\ref{or:HB} gives
\begin{equation}\label{sm:highX-terminal-upper}
 W<\frac{f\,2(25/48)95\cdot102/(6250-190)+\rho_U}{1-B_{95}}.
\end{equation}

Put $z_\pm=u_\pm-\beta$, $R=17/98$, and $m=19/490$.
The pair satisfies
$|z_\pm|<R$, $|(z_++z_-)/2|<m$, and $\Re z_\pm\le0$.
Its insets lie in $[3/4,6/5]$: for the lower bound use
\[
 I(\beta+z)\ge(1-\rho_U)(1-R)-\rho_UR^2/2>3/4,
\]
and for the upper bound use $I(\beta+z)\le1+R<6/5$.
On this whole interval,
$\sqrt I\ge I/2+49/100$.
Indeed both sides are nonnegative, and the difference
$(I/2+49/100)^2-I$ is a convex quadratic negative at
the two endpoints.
The bound on the pair mean gives
\begin{align*}
 I_++I_-&>I_\Sigma:=2(1-\rho_U)(1-m)-\rho_UR^2,\\
 I_\Sigma&=\frac{43251409}{24010000}>0,\\
 |s-u_+|^2+|s-u_-|^2
 &<D_\Sigma:=2(1+\delta^2)+4(7/6)m+2R^2,
\end{align*}
where $(7/6)^2>1+\delta^2$.
Since $I_\Sigma/2+49/50>0$, Cauchy--Schwarz on these two terms,
together with the separate source $\beta$, yields
\begin{equation}\label{sm:highX-terminal-lower}
 E>\frac{1-\rho_U}{1+\delta^2}
                    +\frac{(I_\Sigma/2+49/50)^2}{D_\Sigma}.
\end{equation}
The exact lower bound in \eqref{sm:highX-terminal-lower}
exceeds the upper bound in \eqref{sm:highX-terminal-upper} by
\[
 \frac{9014845822117732608757183663183337}
      {97049415152918630536485588216000000}>\frac{92}{1000}.
\]
This contradicts $W\ge E$.

The source bound, the polynomial
\eqref{sm:highX-q-polynomial}, the geometric and phase
bounds, and the two terminal bounds are rationally checked
by \cert{certify_T6251_highX_bootstrap.py}; its code is also
called by the complete ten-cell verifier.
For $N\ge6250$, the initial fields and the source strengthen,
while the finite polar and radial corrections decrease.
On the terminal range $h<95$, the denominator satisfies
$N-2h>N-190\ge6060$. All other
domains and contours are fixed, and every scale and moment
endpoint has been included. This proves the lemma on the
entire stated degree ray.
\end{proof}
